\documentclass[12pt,a4paper]{article}
\usepackage[T1]{fontenc}
\usepackage{amsmath,amssymb,mathrsfs,tikz,times,pifont}
\usepackage{enumitem}
\usepackage{multicol}
\usepackage{lmodern}
\usetikzlibrary{trees}
\newcommand\circitem[1]{%
\tikz[baseline=(char.base)]{
\node[circle,draw=gray, fill=red!55,
minimum size=1.2em,inner sep=0] (char) {#1};}}
\newcommand\boxitem[1]{%
\tikz[baseline=(char.base)]{
\node[fill=cyan,
minimum size=1.2em,inner sep=0] (char) {#1};}}
\setlist[enumerate,1]{label=\protect\circitem{\arabic*}}
\setlist[enumerate,2]{label=\protect\boxitem{\alph*}}
%%%::::::by chnini ameur :::::::%%%
\everymath{\displaystyle}
\usepackage[left=1cm,right=1cm,top=1cm,bottom=1.7cm]{geometry}
\usepackage[colorlinks=true, linkcolor=blue, urlcolor=blue, citecolor=blue]{hyperref}
\usepackage{array,multirow}
\usepackage[most]{tcolorbox}
\usepackage{varwidth}
\usepackage{float} %pour utiliser l'option [H] qui force l'image à apparaître exactement à l'endroit où elle est placée dans le code.
\tcbuselibrary{skins,hooks}
\usetikzlibrary{patterns}
%%%::::::by chnini ameur :::::::%%%
\newtcolorbox{exa}[2][]{enhanced,breakable,before skip=2mm,after skip=5mm,
colback=yellow!20!white,colframe=black!20!blue,boxrule=0.5mm,
attach boxed title to top left ={xshift=0.6cm,yshift*=1mm-\tcboxedtitleheight},
fonttitle=\bfseries,
title={#2},#1,
% varwidth boxed title*=-3cm,
boxed title style={frame code={
\path[fill=tcbcolback!30!black]
([yshift=-1mm,xshift=-1mm]frame.north west)
arc[start angle=0,end angle=180,radius=1mm]
([yshift=-1mm,xshift=1mm]frame.north east)
arc[start angle=180,end angle=0,radius=1mm];
\path[left color=tcbcolback!60!black,right color = tcbcolback!60!black,
middle color = tcbcolback!80!black]
([xshift=-2mm]frame.north west) -- ([xshift=2mm]frame.north east)
[rounded corners=1mm]-- ([xshift=1mm,yshift=-1mm]frame.north east)
-- (frame.south east) -- (frame.south west)
-- ([xshift=-1mm,yshift=-1mm]frame.north west)
[sharp corners]-- cycle;
},interior engine=empty,
},interior style={top color=yellow!5}}
%%%%%%%%%%%%%%%%%%%%%%%
\usepackage{fancyhdr}
\usepackage{eso-pic}         % Pour ajouter des éléments en arrière-plan
% Commande pour ajouter du texte en arrière-plan
\usepackage{tkz-tab}
\AddToShipoutPicture{
    \AtTextCenter{%
        \makebox[0pt]{\rotatebox{80}{\textcolor[gray]{0.7}{\fontsize{5cm}{5cm}\selectfont PGB}}}
    }
}
\usepackage{lastpage}
\fancyhf{}
\pagestyle{fancy}
\renewcommand{\footrulewidth}{1pt}
\renewcommand{\headrulewidth}{0pt}
\renewcommand{\footruleskip}{10pt}
\fancyfoot[R]{
\color{blue}\ding{45}\ \textbf{2025}
}
\fancyfoot[L]{
\color{blue}\ding{45}\ \textbf{Prof:M. BA}
}
\cfoot{\bf
\thepage /
\pageref{LastPage}}
% Création du compteur pour les exercices
\newcounter{exercice}
\renewcommand{\theexercice}{\arabic{exercice}}  % Définit l'affichage du compteur en chiffres arabes

% Définir la commande \exo
\newcommand{\exo}{\refstepcounter{exercice}\textbf{Exercice \theexercice} }
\begin{document}
\renewcommand{\arraystretch}{1.5}
\renewcommand{\arrayrulewidth}{1.2pt}
\begin{tikzpicture}[overlay,remember picture]
    \node[draw=blue,line width=1.2pt,fill=purple,text=blue,inner sep=3mm,rounded corners,pattern=dots]at ([yshift=-2.5cm]current page.north) {\begingroup\setlength{\fboxsep}{0pt}\colorbox{white}{\begin{tabular}{|*1{>{\centering \arraybackslash}p{0.28\textwidth}} |*2{>{\centering \arraybackslash}p{0.2\textwidth}|} *1{>{\centering \arraybackslash}p{0.19\textwidth}|} }
                \hline
                \multicolumn{3}{|c|}{$\diamond$$\diamond$$\diamond$\ \textbf{Lycée de Dindéfélo}\ $\diamond$$\diamond$$\diamond$ } & \textbf{A.S. : 2026/2027}                                              \\ \hline
                \textbf{Matière: Mathématiques}                                                                                    & \textbf{Niveau : 1}\textbf{L} & \textbf{Date: 05/11/2026} & \textbf{} \\ \hline
                \multicolumn{4}{|c|}{\parbox[c]{10cm}{\begin{center}
                                                                  \textbf{{\Large\sffamily Composition Des Applications}}
                                                              \end{center}}}                                                                                                        \\ \hline
            \end{tabular}}\endgroup};
\end{tikzpicture}

\vspace{3cm}
\fbox{\textbf{\exo}}
Choisir la bonne réponse
\begin{enumerate}
\item Soient $f$ et $g$ deux applications de $I$ définies par $f(x) = -\sqrt{x} \,;\, I = [0 \,;\, +\infty[$ et $g(x) = x^2 + 1 \,;\, I = \mathbb{R}$. La composée $f \circ g(x)$ est égale à \\[0.5em]
$\text{a)}-\sqrt{x^2 + 1} \hspace{1.5cm} \text{b)} \ x^2 + 1 \hspace{1.5cm} \text{c)} \ (-\sqrt{x})^2 - 1$
\item On considère les fonctions $f$ et $g$ telles que $f(x) = x + 1$ et $g(x) = x^2 - 1$ alors on a :

a) $g \circ f(x) = x^2 + 2x \hspace{1cm} \text{b)} \ g \circ f(x) = x^2 \hspace{1cm} \text{c)} \ g \circ f(x) = x^2 + x$
\item On considère les fonctions $f$ et $g$ telles que $f(x) = \frac{1}{x}$ et $g(x) = x - 3$ alors on a :

a) $f \circ g(x) = \frac{1}{x-3} \hspace{1cm} \text{b)} \ f \circ g(x) = \frac{1}{x} - 3 \hspace{1cm} \text{c)} \ f \circ g(x) = \frac{x}{x-3}$
\item Soit $\frac{x^{2}+3x+3}{x+2}$

a) $(f \circ g)(x) = \frac{x^{2}-3x+3}{x+2} \hspace{1cm} \text{b)} \ (f \circ g)(x) = \frac{-x^{2}+3x-3}{x-2} \hspace{1cm} \text{c)} \ (f \circ g)(x) = \frac{x^{2}-3x+3}{x-2}$
\end{enumerate}

\fbox{\textbf{\exo}}

On considère deux fonctions définies par : $f(x) = 2x^2 + 1 \quad \text{et} \quad g(x) = \sqrt{x - 3}$

\begin{enumerate}
    \item Déterminer le domaine de définition de $f$ et de $g$ 
    \item Donner les expressions de : $f \circ g$ et $g \circ f$ 
    \item Calculer $f \circ g(3)$ et $g \circ f(-2)$ 
\end{enumerate}

\fbox{\textbf{\exo}}

On donne les applications $f$ et $g$ définies sur $]2 ; +\infty[$ par :

$f(x) = -x^2 + 3x \quad \text{et} \quad g(x) = \frac{x+1}{x-2}$.

\begin{enumerate}
    \item Calculer $g \circ f(x)$ et $f \circ g(x)$.
    \item Déterminer de deux façons différentes l'image de $-1$ par $g \circ f$.
\end{enumerate}
\end{document}